% Exactly six Lean frames. No footnotes or footer notes.
\begin{frame}{Lean 1/6: from economics to a precise statement}
\begin{center}
\Large Paper claim $\;\longrightarrow\;$ Lean statement $\;\longrightarrow\;$ Proof
\end{center}
\vspace{6mm}
\begin{points}
\item A \texttt{Spec : Prop} states objects, assumptions, and conclusion.
\item Lean checks a proof of that exact formal statement.
\item Matching the statement to the paper requires a separate audit.
\end{points}
\vspace{4mm}
Our run checks supporting algebra and a finite-model threshold.
\end{frame}

\begin{frame}[fragile]{Lean 2/6: the bottom-gain threshold}
\[0<a<\tfrac13\quad\Longrightarrow\quad
\left(\tfrac15<\frac{a}{2(1-a)}\iff\tfrac27<a\right).\]
\lstinputlisting[firstline=21,lastline=30,basicstyle=\ttfamily\footnotesize]{lean/DiscreteCalculations.lean}
\medskip
\texttt{bottomGain} proves both directions over real numbers.\par
\textbf{Build: PASS.} This is the finite wage inequality, not Proposition 5.
\end{frame}

\begin{frame}[fragile]{Lean 3/6: why the inequality proof works}
\lstinputlisting[firstline=26,lastline=30]{lean/DiscreteCalculations.lean}
\vspace{5mm}
\begin{points}
\item \texttt{hden}: prove $2(1-a)>0$ from the stated domain.
\item \texttt{lt\_div\_iff\textsubscript{0}}: clear division without reversing the inequality.
\item \texttt{constructor}: split ``if and only if'' into two directions.\par
\texttt{linarith} closes each direction by linear arithmetic.
\end{points}
\end{frame}

\begin{frame}[fragile]{Lean 4/6: zero profit implies a wage identity}
\[n(z)=\frac1{h(1-z)},\qquad n(z)(a-w)-r=0
\quad\Longrightarrow\quad w=a-h(1-z)r.\]
\lstinputlisting[firstline=28,lastline=32,basicstyle=\ttfamily\footnotesize]{lean/FirmAlgebra.lean}
\medskip
\texttt{zeroProfitWage} uses $h>0$ and $z<1$ to justify division.\par\medskip
\textbf{Build: PASS.} It checks conditional algebra over real scalars;\par
it does not establish market clearing or equilibrium rents.
\end{frame}

\begin{frame}{Lean 5/6: a missing assumption changes the claim}
Dropping $z<1$ allows
\[h=\tfrac12,\quad z=1,\quad a=\tfrac12,\quad w=1,\quad r=0.\]
\begin{points}
\item Lean's real division is total: here $n(z)=0$.
\item Profit is zero, but the proposed wage identity says $1=1/2$.
\item The separate counterexample compiles successfully.
\end{points}
\vspace{4mm}
This tests unrestricted algebra; it is not an admissible economic team.
\end{frame}

\begin{frame}{Lean 6/6: what passed and what remains open}
\begin{center}
\begin{tabular}{lp{8cm}}\toprule
Status & Evidence\\\midrule
PASS & Support proofs, full build, and required fast check.\\
STOPPED & Semantic preflight: missing source statement map.\\
0 of 6 & Full continuum propositions proved.\\\bottomrule
\end{tabular}
\end{center}
\vspace{4mm}
Matching, equilibrium, sorting, and source-level audits remain open.
\begin{block}{What I learned}
A checked proof is useful when its assumptions and scope are explicit.
\end{block}
\end{frame}
